\section{Paper 2: Học tăng cường và ứng dụng trong các lĩnh
vực}\label{paper-2-hux1ecdc-tux103ng-cux1b0ux1eddng-vuxe0-ux1ee9ng-dux1ee5ng-trong-cuxe1c-lux129nh-vux1ef1c}

\subsection{Hồ sơ bài báo và mục đích
đọc}\label{hux1ed3-sux1a1-buxe0i-buxe1o-vuxe0-mux1ee5c-ux111uxedch-ux111ux1ecdc}

Bài \emph{A Gentle Introduction to Reinforcement Learning and Its
Application in Different Fields} do Muddasar Naeem, Syed Tahir Hussain
Rizvi và Antonio Coronato thực hiện, đăng trên IEEE Access, tập 8, năm
2020, trang 209320--209344, DOI 10.1109/ACCESS.2020.3038605. Bản PDF
được cung cấp có 25 trang, gồm nền tảng MDP, nguyên lý RL, các nhóm
thuật toán, ứng dụng và thảo luận. \cite{naeem}

Mục đích của bài là cung cấp một lộ trình nhập môn tương đối rộng, giúp
người đọc hiểu những khái niệm trước khi bước vào thuật toán chuyên sâu.
Đóng góp phù hợp nhất để đánh giá là cách hệ thống hóa và nối kiến thức
với ứng dụng. Bài không cung cấp một benchmark thống nhất chạy mọi thuật
toán trên mọi lĩnh vực, nên không thể suy ra một bảng xếp hạng phổ quát
từ số lượng ví dụ được nhắc đến.

Báo cáo đọc paper theo chuỗi: mô hình quyết định → đại lượng cần học →
cách ước lượng → điều kiện áp dụng → kiểm tra bằng assignment. Các công
thức được viết lại với ký hiệu nhất quán và kỳ vọng rõ ràng. Đây là phần
diễn giải kỹ thuật của báo cáo, giúp tránh những giản lược ký hiệu trong
tài liệu nhập môn. Ví dụ số và kết quả mô phỏng được phân biệt rõ với
nội dung nguyên bản.

\subsection{RL khác học có giám sát và học không giám
sát}\label{rl-khuxe1c-hux1ecdc-cuxf3-giuxe1m-suxe1t-vuxe0-hux1ecdc-khuxf4ng-giuxe1m-suxe1t}

Trong học có giám sát, một quan sát đi cùng nhãn hoặc giá trị mục tiêu.
Mô hình học quan hệ giữa đặc trưng và nhãn trên tập dữ liệu cho trước.
Trong học không giám sát, mục tiêu thường là tìm cấu trúc, biểu diễn
hoặc nhóm trong dữ liệu không có nhãn tác vụ. Trong RL, tác tử chọn hành
động rồi nhận phản hồi; hành động ảnh hưởng dữ liệu quan sát tiếp theo,
vì vậy quá trình thu thập dữ liệu gắn với chính sách hiện tại.

Phần thưởng khác nhãn đáp án. Khi cửa hàng giảm giá và nhận lợi nhuận
30, môi trường không nói hành động còn lại sẽ cho lợi nhuận bao nhiêu
trong cùng tình huống. Tác tử phải học thông qua nhiều lần tương tác,
đối mặt với nhiễu và đánh đổi giữa thử mới với khai thác điều đã biết.
Nếu hậu quả xuất hiện chậm, còn phải xác định hành động nào trong chuỗi
đóng góp cho kết quả cuối.

Tutorial intent sử dụng \textbf{supervised learning}: mỗi câu có nhãn
ecommerce, tourism hoặc healthcare administration và mạng tối ưu cross
entropy. Assignment sử dụng \textbf{reinforcement learning}: agent chọn
hành động giảm giá và học từ lợi nhuận ngẫu nhiên. Hai phần đều có quá
trình training nhưng không nên gọi classifier là RL. Tương tự, một hệ có
nhiều agent không tự động là MARL nếu các thành viên không cập nhật
chính sách từ tương tác.

\subsection{Mô hình MDP và giả định
Markov}\label{muxf4-huxecnh-mdp-vuxe0-giux1ea3-ux111ux1ecbnh-markov}

Một MDP thường được biểu diễn bằng
\((\mathcal S,\mathcal A,P,R,\gamma,\rho_0)\): không gian trạng thái,
hành động, xác suất chuyển, phần thưởng kỳ vọng, hệ số chiết khấu và
phân phối khởi đầu. Với bài hữu hạn chân trời, cần nêu thêm \(H\) hoặc
dùng trạng thái chứa thời gian. MDP là mô hình bài toán; RL là một họ
phương pháp học quyết định, thường khi không biết đầy đủ \(P\) và \(R\).

Tính Markov có thể viết:

\[P(S_{t+1},R_{t+1}\mid S_0,A_0,\ldots,S_t,A_t)
=P(S_{t+1},R_{t+1}\mid S_t,A_t).\]

Điều kiện không nói quá khứ không quan trọng trong thế giới thực, mà nói
trạng thái đã tóm tắt đủ thông tin quá khứ liên quan đến tương lai. Nếu
trạng thái thiếu tồn kho, đơn hàng đang xử lý hoặc hành động trước mà
đối thủ dùng để phản ứng, điều kiện có thể không còn. Một chuỗi quan sát
có dạng vector không mặc nhiên là trạng thái Markov.

Cần tách tính Markov và tính dừng. Markov liên quan đến thông tin cần
điều kiện hóa; tính dừng liên quan đến việc cùng một trạng thái và hành
động có cùng quy luật qua thời gian. Một quá trình có thể Markov nhưng
quy luật phụ thuộc thời điểm. Có thể bổ sung thời gian hoặc chế độ vào
trạng thái để mô hình hóa sự phụ thuộc đó, nhưng khi không quan sát chế
độ thì bài toán vẫn có bất định.

\subsection{Những thành phần cơ bản của môi
trường}\label{nhux1eefng-thuxe0nh-phux1ea7n-cux1a1-bux1ea3n-cux1ee7a-muxf4i-trux1b0ux1eddng}

\textbf{Trạng thái} phải chứa thông tin cần cho quyết định, không nhất
thiết là mọi dữ liệu có thể lưu. Thêm quá nhiều biến làm tăng độ phức
tạp; bỏ biến quan trọng làm giá trị ước lượng bị trộn. Trong assignment,
quan sát tham chiếu có hai bit nên chỉ có bốn trạng thái cục bộ. Khi
kiểm tra giả định Markov, báo cáo chỉ ra cần thêm hành động trước của
chính mình để mô tả phản ứng của rule opponent.

\textbf{Hành động} là lựa chọn thực sự được môi trường cho phép. Ở đây
\(a\in\{0,1\}\), không giảm và giảm 20\%. Nếu mở rộng sang nhiều mức
giảm, cần kiểm tra biên lợi nhuận và tập hợp hành động hợp lệ. Một mạng
xuất ra giá trị liên tục không tự giải quyết tính hợp lệ của giao dịch
hoặc hạn mức kinh doanh.

\textbf{Chuyển trạng thái} mô tả điều xảy ra sau hành động. Nhu cầu
trong mô phỏng được lấy độc lập với hành động, nhưng hành động hiện tại
trở thành lịch sử quan sát ở bước sau. Vì rule opponent phản ứng theo
lịch sử này, agent vẫn có ảnh hưởng đến tương lai dù không điều khiển
nhu cầu. Chỉ nhìn nhu cầu độc lập rồi kết luận bài toán là bandit sẽ bỏ
qua cơ chế phản hồi đối thủ.

\textbf{Phần thưởng} phản ánh mục tiêu đo được. Trong mô phỏng là lợi
nhuận số học, không phải doanh thu hoặc số đơn. Việc chọn reward là lựa
chọn mô hình hóa có hệ quả: tối ưu số lượng bán sẽ có động lực khác tối
ưu lợi nhuận; thêm phạt tồn kho hay phạt biến động giá lại tạo chính
sách khác. Cần xác định mục tiêu trước khi chọn thuật toán.

\subsection{Return, hệ số chiết khấu và tiêu chí kinh
doanh}\label{return-hux1ec7-sux1ed1-chiux1ebft-khux1ea5u-vuxe0-tiuxeau-chuxed-kinh-doanh}

Return hữu hạn bắt đầu từ bước \(t\) được viết:

\[G_t=\sum_{k=0}^{H-t-1}\gamma^k R_{t+k+1}.\]

Khi \(\gamma=0\), chỉ phần thưởng kế tiếp được xem xét. Khi
\(0<\gamma<1\), các phần thưởng xa hơn có trọng số nhỏ hơn và chuỗi vô
hạn bị chặn nếu reward bị chặn. Khi \(\gamma=1\), nhiệm vụ hữu hạn vẫn
hợp lệ; trong nhiệm vụ vô hạn phải có điều kiện bổ sung để tổng có
nghĩa. Vì vậy không thể nói mọi lựa chọn \(\gamma=1\) đều sai hoặc đều
an toàn.

Có một khác biệt quan trọng trong assignment: cập nhật Q dùng
\(\gamma=0.9\), trong khi báo cáo hiệu năng yêu cầu lợi nhuận trung bình
không chiết khấu qua 20 bước. Hai tiêu chí có liên quan nhưng không đồng
nhất trong mọi môi trường. Nếu một hành động hy sinh lợi nhuận đầu kỳ để
có lợi ích cuối kỳ, việc giảm trọng số tương lai có thể thay đổi lựa
chọn tối ưu theo chỉ số kinh doanh.

Trong mô hình hiện tại, quy hoạch động xác nhận chính sách không giảm
giá là best response với rule opponent ở các tham số được dùng. Do đó
khác biệt tiêu chí không làm thay đổi mốc hành động này. Tuy nhiên, báo
cáo vẫn nêu sự không đồng nhất vì nó quan trọng khi thêm tồn kho, ngân
sách khuyến mãi hoặc khách hàng quay lại.

Ví dụ của báo cáo: chuỗi reward \((5,5,20)\) có tổng 30, nhưng return
với \(\gamma=0.9\) là \(5+0.9\times5+0.9^2\times20=25.7\). Một chính
sách khác nhận \((12,10,6)\) có tổng 28 và return 25.86. Như vậy chính
sách có tổng ít hơn vẫn có thể được ưu tiên dưới chiết khấu. Ví dụ này
cho thấy phải công bố rõ tối ưu đại lượng nào.

\subsection{Chính sách, kế hoạch và hàm giá
trị}\label{chuxednh-suxe1ch-kux1ebf-houx1ea1ch-vuxe0-huxe0m-giuxe1-trux1ecb}

Chính sách xác định cách hành động theo trạng thái hoặc lịch sử. Chính
sách xác định viết \(a=\pi(s)\); chính sách ngẫu nhiên viết
\(\pi(a\mid s)\). Một kế hoạch là chuỗi hoặc đồ thị bước để thực hiện,
còn chính sách là quy tắc phản ứng cho những tình huống có thể gặp.
Trong môi trường ngẫu nhiên, một danh sách hành động cố định thường
không đủ thay thế chính sách.

Hàm giá trị trạng thái và hành động được định nghĩa bằng kỳ vọng:

\[V^\pi(s)=\mathbb E_\pi[G_t\mid S_t=s],\qquad
Q^\pi(s,a)=\mathbb E_\pi[G_t\mid S_t=s,A_t=a].\]

Một reward lớn trong một lượt không bằng Q lớn. Q phản ánh lợi ích dài
hạn kỳ vọng khi thực hiện hành động rồi tiếp tục theo chính sách. Mẫu
sales có nhiễu, nên hai lần cùng trạng thái và hành động vẫn có thể cho
reward khác nhau. Bảng Q giảm bớt nhiễu nhờ cập nhật qua nhiều quan sát,
nhưng cách giảm phụ thuộc learning rate và cấu trúc bài toán.

Hàm giá trị tối ưu \(V^*\), \(Q^*\) dùng chính sách tốt nhất trong lớp
đang xét và giả định mô hình. Nếu trạng thái thiếu thông tin, một chính
sách tối ưu trên biểu diễn đó có thể kém chính sách dùng bộ nhớ đầy đủ.
Từ ``tối ưu'' luôn cần đi kèm điều kiện: tối ưu đối với trạng thái nào,
môi trường nào, đối thủ nào và mục tiêu nào.

\subsection{Phương trình Bellman}\label{phux1b0ux1a1ng-truxecnh-bellman}

Với reward kỳ vọng \(R(s,a)\) và chuyển trạng thái \(P(s'\mid s,a)\),
phương trình kỳ vọng Bellman có dạng:

\[V^\pi(s)=\sum_a\pi(a\mid s)\left[R(s,a)+\gamma\sum_{s'}P(s'\mid s,a)V^\pi(s')\right].\]

Phương trình tối ưu thay trung bình theo chính sách bằng phép cực đại:

\[V^*(s)=\max_a\left[R(s,a)+\gamma\sum_{s'}P(s'\mid s,a)V^*(s')\right].\]

Đối với Q, ta có:

\[Q^*(s,a)=R(s,a)+\gamma\sum_{s'}P(s'\mid s,a)\max_{a'}Q^*(s',a').\]

Các biểu thức này làm rõ hai thành phần: lợi ích tức thời và giá trị của
trạng thái tiếp theo. Trong một số phần nhập môn của paper, ký hiệu được
rút gọn để dễ trình bày. Báo cáo thêm tổng xác suất và điều kiện chính
sách để tránh hiểu rằng chỉ cộng mọi giá trị tương lai với trọng số bằng
nhau.

Với chân trời hữu hạn, nên viết \(V_t\) thay vì một \(V\) không phụ
thuộc thời gian, và điều kiện biên \(V_H=0\). Khi cùng một quan sát xuất
hiện ở đầu và cuối episode, giá trị phần thưởng còn có thể nhận khác
nhau. Một bảng Q không chứa thời gian nhưng nhận cả target terminal lẫn
nonterminal sẽ trộn các ngữ cảnh này; đây là hạn chế của baseline được
giữ lại để phân tích.

\subsection{Dynamic programming: biết mô hình để lập kế
hoạch}\label{dynamic-programming-biux1ebft-muxf4-huxecnh-ux111ux1ec3-lux1eadp-kux1ebf-houx1ea1ch}

Section IV-A trình bày dynamic programming, gồm đánh giá chính sách, cải
thiện chính sách, policy iteration và value iteration. Khi biết mô hình,
ta có thể tính kỳ vọng thay vì chờ mẫu thực nghiệm. Policy evaluation
giữ nguyên \(\pi\) và giải giá trị; policy improvement chọn hành động
tốt hơn theo giá trị hiện có; policy iteration luân phiên hai bước.

Value iteration áp dụng trực tiếp toán tử tối ưu Bellman. Trong MDP hữu
hạn có chiết khấu thích hợp, phép lặp có nền tảng hội tụ. Tuy nhiên, khi
mô hình lớn hoặc không biết, liệt kê toàn bộ trạng thái và chuyển tiếp
có thể không khả thi. Đây là lý do học từ mẫu trở nên cần thiết, chứ
không phải vì quy hoạch động vốn kém thông minh hơn RL.

Trong assignment, mô hình chỉ có hai mức nhu cầu, hai hành động và phản
ứng rule rõ ràng. Vì thế quy hoạch động hữu hạn chân trời là mốc kiểm
tra rất mạnh. Nó giúp phân biệt việc ``Q-learning đã học được một chiến
lược'' với việc ``chiến lược đó đạt best response''. Một thuật toán học
phức tạp mà không so với mốc đơn giản có thể tạo ra kết luận sai về hiệu
quả.

Mã DP của báo cáo dùng trạng thái \((d,a_{t-1}^{A},a_{t-1}^{B})\) và
thời điểm \(t\). Với rule opponent, hành động hiện tại
\(b=d\,a_{t-1}^{A}\). Với mỗi lựa chọn \(a\), thuật toán tính reward kỳ
vọng và trung bình giá trị ở hai mức nhu cầu kế tiếp. Việc thực hiện từ
\(t=H-1\) về 0 cho ra chính sách phụ thuộc thời gian một cách chính xác
trong mô hình.

\subsection{Monte Carlo: học từ return hoàn
chỉnh}\label{monte-carlo-hux1ecdc-tux1eeb-return-houxe0n-chux1ec9nh}

Section IV-B giới thiệu Monte Carlo như cách ước lượng giá trị qua các
episode hoàn tất. Với mỗi lần thăm trạng thái, tính return thực tế rồi
lấy trung bình. Nếu dữ liệu được lấy đúng từ chính sách mục tiêu và
episode phù hợp, trung bình mẫu là cách tự nhiên để ước lượng kỳ vọng.
Không cần biết bảng xác suất chuyển trạng thái.

Ưu điểm là không bootstrap từ một giá trị ước lượng khác. Nhược điểm là
phải chờ tương lai của episode và return có thể biến động mạnh theo
nhiều bước ngẫu nhiên. Nếu một episode dài gồm nhiều thành phần nhiễu,
mẫu return ở đầu kỳ có phương sai lớn. First-visit và every-visit phân
biệt cách dùng những lần gặp lại trạng thái trong một episode.

Off-policy Monte Carlo cần xử lý khác biệt giữa chính sách thu thập và
chính sách đánh giá, thường qua importance sampling với điều kiện hỗ trợ
thích hợp. Trọng số có thể rất lớn khi chuỗi dài hoặc chính sách khác
nhau mạnh. Do đó việc có log lịch sử không tự động cho phép đánh giá mọi
chính sách mới một cách ổn định.

Liên hệ bài lab: mỗi episode có 20 bước nên có thể bổ sung MC để đối
chiếu với Q-learning. Tuy vậy, bản thực nghiệm chính không chạy MC; báo
cáo không gán cho nó số liệu chưa đo. Đề xuất có thể thực hiện là giữ
cùng seed và budget tương tác, thêm MC control với cùng không gian quan
sát rồi so mức sai lệch so với oracle.

\subsection{Temporal difference và ý nghĩa
bootstrapping}\label{temporal-difference-vuxe0-uxfd-nghux129a-bootstrapping}

Section IV-C trình bày temporal difference, kết hợp việc học từ mẫu với
ước lượng giá trị tương lai. TD(0) cho trạng thái có dạng:

\[V(S_t)\leftarrow V(S_t)+\alpha\left[R_{t+1}+\gamma V(S_{t+1})-V(S_t)\right].\]

Đại lượng trong ngoặc là TD error. Khi target cao hơn dự đoán, giá trị
tăng; khi thấp hơn, giá trị giảm. Không cần chờ episode hoàn tất, nhưng
target phụ thuộc ước lượng hiện tại nên có thể mang bias. Nhận định ``TD
phương sai thấp hơn MC'' là trực giác thường hữu ích, không phải bất
đẳng thức phổ quát cho mọi cấu hình dữ liệu và estimator.

Bootstrapping là dùng một giá trị ước lượng để cập nhật giá trị khác.
Trong Q-learning, target có phép max trên trạng thái kế tiếp. Nó truyền
thông tin từ tương lai về các hành động trước. Nhưng nếu ước lượng kế
tiếp không phù hợp do trạng thái thiếu hoặc môi trường thay đổi, sai
lệch cũng có thể truyền ngược.

Ở terminal thật, giá trị tương lai bằng 0. Trong nhiệm vụ continuing bị
cắt chỉ để quản lý thời lượng, ``hết 20 bước'' có thể là truncation và
việc bootstrap cần khác terminal. Đề tham chiếu dùng episode hữu hạn với
terminal ở bước cuối, nên báo cáo giữ quy ước này; đồng thời chỉ ra cần
thêm thời gian vào trạng thái nếu muốn biểu diễn chính xác nhiệm vụ hữu
hạn.

\subsection{SARSA và Q-learning}\label{sarsa-vuxe0-q-learning}

SARSA cập nhật theo hành động thực sự được chọn ở trạng thái tiếp theo:

\[Q(S_t,A_t)\leftarrow Q(S_t,A_t)+\alpha[R_{t+1}+\gamma Q(S_{t+1},A_{t+1})-Q(S_t,A_t)].\]

Nó gắn với chính sách đang hành động, gồm cả khám phá, nên thường được
gọi là on-policy. Q-learning thay \(Q(S_{t+1},A_{t+1})\) bằng
\(\max_aQ(S_{t+1},a)\), hướng đến chính sách greedy dù dữ liệu có thể
đến từ epsilon-greedy, nên là off-policy trong cách phân loại này.

Sự khác nhau không nói SARSA luôn an toàn còn Q-learning luôn rủi ro.
Kết quả phụ thuộc môi trường, reward, biểu diễn và quy trình đánh giá.
Ví dụ kinh điển về hành động khám phá nguy hiểm chỉ minh họa một cấu
hình cụ thể. Khi thảo luận một ứng dụng kinh doanh, cần đo chi phí khám
phá và ràng buộc thực tế.

Ví dụ cập nhật theo đề: \(Q=4\), \(\alpha=0.1\), reward 6,
\(\gamma=0.9\), giá trị tốt nhất kế tiếp 5. Target là
\(6+0.9\times5=10.5\); TD error là 6.5; giá trị mới
\(4+0.1\times6.5=4.65\). Nếu cùng transition kết thúc thật, target chỉ
là 6 và giá trị mới 4.2. Hai ca được kiểm thử riêng trong mã.

\subsection{Epsilon-greedy và khám
phá}\label{epsilon-greedy-vuxe0-khuxe1m-phuxe1}

Với hai hành động, epsilon-greedy chọn một hành động ngẫu nhiên trong
hai lựa chọn với xác suất \(\epsilon\), còn lại chọn hành động tốt nhất
theo bảng. Vì nhánh ngẫu nhiên vẫn có thể chọn hành động greedy, xác
suất chọn hành động không greedy là \(\epsilon/2\) khi chỉ có một greedy
action. Nói epsilon bằng đúng tỷ lệ hành động sai sẽ không chính xác.

Lịch tham chiếu giảm theo episode:
\(\epsilon_e=\max(0.02,0.2\times0.995^e)\) với \(e\) bắt đầu từ 0. Giải
\(0.2\times0.995^e\le0.02\) cho thấy ngưỡng sàn đạt xấp xỉ episode chỉ
số 460. Hơn 4.500 episode còn lại vẫn giữ khám phá nhỏ. Khi đánh giá,
báo cáo đặt epsilon bằng 0 và không cập nhật Q.

Khám phá quá ít có thể làm bỏ sót hành động, nhất là trạng thái hiếm.
Khám phá quá nhiều gây mất reward trong training và có thể ảnh hưởng
hành vi đối thủ. Trong bài bán hàng, một hành động giảm giá hiện tại còn
làm rule opponent có khả năng giảm ở bước sau khi nhu cầu cao. Vì vậy
chi phí khám phá không chỉ là margin của bước hiện tại.

Báo cáo chạy ba cấu hình khám phá: lịch giảm tham chiếu, epsilon cố định
0.05 và 0.30. Việc giữ cố định learning rate, số episode và các seed
giúp so sánh có ý nghĩa. Dù vậy, một chênh lệch sau năm seed chưa đủ
chứng minh cơ chế khám phá là nguyên nhân duy nhất; kết quả còn chịu
nhiễu Q, aliasing và trộn thời điểm terminal.

\subsection{Learning rate và điều kiện hội
tụ}\label{learning-rate-vuxe0-ux111iux1ec1u-kiux1ec7n-hux1ed9i-tux1ee5}

Learning rate điều khiển mức thay đổi trước một mẫu mới. Với target cố
định \(y\), cập nhật \(Q\leftarrow(1-\alpha)Q+\alpha y\) là một trung
bình có trọng số. Alpha lớn phản ứng mạnh, alpha nhỏ tích lũy chậm.
Trong thực tế target cũng thay đổi theo Q và dữ liệu nên trực giác trung
bình chỉ giải thích một phần hành vi.

Các bảo đảm hội tụ cổ điển cho tabular Q-learning đòi hỏi những điều
kiện về môi trường Markov và dừng, độ bao phủ state--action và lịch
learning rate thỏa các tổng thích hợp, thường \(\sum\alpha_t=\infty\),
\(\sum\alpha_t^2<\infty\). Alpha cố định 0.1 và số episode hữu hạn của
lab không đủ để viện dẫn trực tiếp bảo đảm hội tụ về \(Q^*\).
\cite{watkins}

Trong lab, alpha 0.01 có thể cần nhiều mẫu hơn để phản ánh reward; alpha
0.50 có thể nhạy với nhiễu sales và các target terminal thấp. Nhìn một
đường reward được làm trơn chưa xác định được Q đã ổn định. Nên kết hợp
đánh giá đóng băng, phân phối chính sách qua seed và độ lệch với oracle.

Cũng cần phân biệt ổn định và tối ưu. Một policy luôn giảm giá có thể
rất ổn định nhưng lợi nhuận thấp. Một policy tốt có thể cho reward thô
dao động lớn vì nhu cầu ngẫu nhiên. Chỉ số thay đổi tham số, tỷ lệ thay
đổi hành động và lợi nhuận đánh giá giải đáp những câu hỏi khác nhau.

\subsection{Actor--critic và policy
gradient}\label{actorcritic-vuxe0-policy-gradient}

Section IV-D giới thiệu actor--critic, nơi actor biểu diễn chính sách
còn critic ước lượng chất lượng để hướng dẫn actor. Một dạng mục tiêu là
\(J(\theta)=\mathbb E_{\pi_\theta}[G_0]\). Gradient có thể được ước
lượng bằng log xác suất hành động nhân một tín hiệu lợi thế. Critic giúp
giảm phương sai hoặc cung cấp tín hiệu thường xuyên hơn.

Actor--critic hữu ích khi hành động liên tục hoặc không gian quá lớn cho
bảng Q. Tuy nhiên, nó thêm mô hình, learning rate và quan hệ phụ thuộc
giữa hai quá trình học. Critic sai có thể hướng actor sai; actor thay
đổi làm dữ liệu critic học thay đổi. Đây không phải phép thay thế mặc
nhiên tốt hơn cho bài có tám giá trị Q.

Một số câu trong phần thảo luận của paper nói về hội tụ ở mức tổng quát.
Báo cáo đọc thận trọng: policy gradient không được bảo đảm tìm global
optimum cho mọi bài toán mạng nơ-ron và mọi ước lượng gradient;
model-based cũng không được bảo đảm đúng mô hình khi lớp mô hình sai.
Cần đặt các phát biểu về hội tụ trong giả định toán học cụ thể.

Liên hệ cải tiến: nếu mở rộng hành động thành mức chiết khấu liên tục và
thêm tồn kho, có thể cân nhắc actor--critic hoặc phương pháp điều khiển
tối ưu. Trước đó cần một môi trường được hiệu chỉnh và baseline đủ mạnh.
Dùng mạng sâu sớm trong môi trường toy làm tăng khó khăn giải thích mà
chưa có lý do về biểu diễn.

\subsection{Bayesian methods và bất
định}\label{bayesian-methods-vuxe0-bux1ea5t-ux111ux1ecbnh}

Section IV-E đề cập phương pháp Bayesian, dùng phân phối niềm tin để
biểu diễn điều chưa biết. Thay vì giữ một con số reward trung bình duy
nhất, có thể giữ posterior cho tham số xác suất chuyển đổi rồi chọn hành
động theo bất định. Thompson sampling là một ví dụ về lấy mẫu giả thuyết
để cân bằng khám phá và khai thác trong các bài phù hợp.

Cần tách \textbf{aleatoric uncertainty} do tính ngẫu nhiên của môi
trường và \textbf{epistemic uncertainty} do thiếu hiểu biết. Dù biết xác
suất mua chính xác, sales binomial vẫn ngẫu nhiên. Thu thập thêm dữ liệu
giúp giảm bất định tham số nhưng không xóa tính ngẫu nhiên của từng lượt
mua. Nhầm hai loại dễ dẫn đến yêu cầu một policy phải thắng trong mọi
episode.

Trong bài bán hàng, nếu mô hình conversion chưa biết, có thể dùng
posterior cho từng ngữ cảnh và hành động. Nhưng đối thủ phản ứng qua
lịch sử khiến mô hình bandit độc lập có thể thiếu động lực dài hạn. Một
cải tiến cần xác định liệu đang học reward tức thời hay toàn bộ quá
trình quyết định có trạng thái.

\subsection{Deep Q-Network và các thành phần ổn
định}\label{deep-q-network-vuxe0-cuxe1c-thuxe0nh-phux1ea7n-ux1ed5n-ux111ux1ecbnh}

Section IV-F trình bày DQN như cách dùng mạng nơ-ron xấp xỉ Q cho không
gian trạng thái lớn. Hai thành phần nổi bật là experience replay và
target network. Replay tái sử dụng transition và giảm tương quan theo
thời gian trong minibatch; target network làm mục tiêu thay đổi chậm hơn
so với việc dùng cùng một mạng cập nhật ngay lập tức.

Một target điển hình có dạng
\(y=r+\gamma\max_{a'}Q_{\bar\theta}(s',a')\) cho transition không
terminal, rồi tối ưu bình phương sai số giữa \(y\) và \(Q_\theta(s,a)\).
DQN vẫn chịu vấn đề xấp xỉ hàm, tối ưu không lồi và dữ liệu phụ thuộc
chính sách. Thêm replay không tự khắc phục non-stationarity khi nhiều
agent cùng học.

Với tám giá trị Q trong lab, mạng sâu sẽ có nhiều tham số hơn đại lượng
cần học và tốn thêm tuning. Bảng Q là lựa chọn hợp lý để người học nhìn
rõ từng giá trị, từng trạng thái và từng lỗi. Bài học từ paper là chọn
công cụ theo kích thước và đặc tính bài toán, không chọn thuật toán mới
chỉ vì có vẻ mạnh.

\subsection{So sánh các họ phương
pháp}\label{so-suxe1nh-cuxe1c-hux1ecd-phux1b0ux1a1ng-phuxe1p}

{\def\LTcaptype{none} % do not increment counter
\begin{longtable}[]{@{}
  >{\raggedright\arraybackslash}p{(\linewidth - 8\tabcolsep) * \real{0.2000}}
  >{\raggedright\arraybackslash}p{(\linewidth - 8\tabcolsep) * \real{0.2000}}
  >{\raggedright\arraybackslash}p{(\linewidth - 8\tabcolsep) * \real{0.2000}}
  >{\raggedright\arraybackslash}p{(\linewidth - 8\tabcolsep) * \real{0.2000}}
  >{\raggedright\arraybackslash}p{(\linewidth - 8\tabcolsep) * \real{0.2000}}@{}}
\toprule\noalign{}
\begin{minipage}[b]{\linewidth}\raggedright
Phương pháp
\end{minipage} & \begin{minipage}[b]{\linewidth}\raggedright
Cần mô hình chuyển tiếp?
\end{minipage} & \begin{minipage}[b]{\linewidth}\raggedright
Target chính
\end{minipage} & \begin{minipage}[b]{\linewidth}\raggedright
Điểm mạnh
\end{minipage} & \begin{minipage}[b]{\linewidth}\raggedright
Hạn chế chính
\end{minipage} \\
\midrule\noalign{}
\endhead
\bottomrule\noalign{}
\endlastfoot
Dynamic programming & Có & Kỳ vọng Bellman & Mốc chính xác khi mô hình
nhỏ, đúng & Chi phí không gian trạng thái, sai mô hình \\
Monte Carlo & Không & Return sau episode & Không bootstrap, diễn giải rõ
& Chờ cuối kỳ, phương sai return \\
TD/SARSA & Không & Reward và giá trị hành động kế tiếp & Cập nhật
online, theo policy hành động & Bootstrap và phụ thuộc biểu diễn \\
Q-learning & Không & Reward và max Q kế tiếp & Tabular nhỏ dễ triển khai
& Không bảo đảm ở mọi môi trường, nhạy trạng thái thiếu \\
Actor--critic & Không nhất thiết & Gradient policy và critic & Phù hợp
policy tham số, hành động lớn & Nhiều nguồn bất ổn và tham số \\
Bayesian & Tùy thiết kế & Posterior và quyết định dưới bất định & Biểu
diễn điều chưa biết & Chi phí suy luận, giả định prior/model \\
DQN & Không cần mô hình tường minh & Bellman target từ mạng mục tiêu &
Xấp xỉ Q trên đầu vào lớn & Data hungry, khó giải thích, cần ổn định \\
\end{longtable}
}

Các trục phân loại không hoàn toàn loại trừ nhau. Bayesian có thể dùng
cho model-based hoặc một thành phần reward; actor--critic có thể dùng mô
hình hoặc không; deep RL có thể value-based hoặc policy-based. Do vậy
không nên vẽ một cây phân loại duy nhất rồi coi mỗi thuật toán chỉ được
đặt vào đúng một lá.

\subsection{Ứng dụng trong paper: đọc bằng cấu trúc bài
toán}\label{ux1ee9ng-dux1ee5ng-trong-paper-ux111ux1ecdc-bux1eb1ng-cux1ea5u-truxfac-buxe0i-touxe1n}

\subsubsection{Healthcare}\label{healthcare}

Paper xem xét lựa chọn điều trị tuần tự, cá nhân hóa và một số quy trình
hỗ trợ. Điểm phù hợp với RL là quyết định hiện tại ảnh hưởng sức khỏe và
lựa chọn ở tương lai. Nhưng việc quan sát dữ liệu điều trị lịch sử không
tương đương có quyền khám phá trực tiếp trên bệnh nhân. Nhiễu quan sát,
confounding và thiếu độ bao phủ hành động khiến đánh giá chính sách khó.

Báo cáo chỉ dùng miền healthcare ở \textbf{phân loại yêu cầu hành
chính}, không tạo chính sách điều trị, chẩn đoán hay khuyến nghị thuốc.
Liên hệ này giúp nhận ra khác biệt mức rủi ro giữa một ví dụ nhập môn và
ứng dụng clinical mà paper khảo sát. Các chi tiết y khoa thực hành nằm
ngoài phạm vi bộ bài.

\subsubsection{Robotics và điều khiển tự
động}\label{robotics-vuxe0-ux111iux1ec1u-khiux1ec3n-tux1ef1-ux111ux1ed9ng}

Robot có chuỗi quan sát--hành động liên tục và hậu quả động học, nên RL
là khung tự nhiên. Phần thưởng có thể gắn với hoàn thành nhiệm vụ, tiêu
hao năng lượng hoặc tránh va chạm. Khó khăn là tương tác thật tốn kém và
hành động sai có thể gây hỏng vật lý. Mô phỏng giúp thu thập dữ liệu
nhưng có khoảng cách sim-to-real.

Một thí nghiệm robot tốt không chỉ báo reward trung bình; cần tỷ lệ
thành công, vi phạm giới hạn, độ bền qua điều kiện tải và khả năng khôi
phục. Các chỉ số này có thể được chuyển thành tư duy cho tutorial: một
hệ trả lời nhanh nhưng không giữ bất biến ngân sách vẫn thất bại nhiệm
vụ.

\subsubsection{Communication và
networking}\label{communication-vuxe0-networking}

Ứng dụng gồm điều khiển tài nguyên, lập lịch, đường truyền hoặc cấu
hình. Trạng thái mạng biến động và quyết định cục bộ có thể ảnh hưởng
người dùng khác. RL phù hợp khi có tương tác lặp và feedback, nhưng
không thể bỏ qua độ trễ quyết định và sự thay đổi phân phối lưu lượng.

Nếu nhiều node học độc lập, throughput của một node phụ thuộc chính sách
của các node còn lại. Khi mọi node cập nhật, mỗi node quan sát một môi
trường hiệu dụng đang đổi. Đây là cầu nối trực tiếp đến vấn đề
non-stationarity trong câu 8 của assignment.

\subsubsection{NLP}\label{nlp}

Paper thảo luận ứng dụng RL trong các tác vụ ngôn ngữ, đối thoại và tạo
chuỗi, nơi đầu ra dài và tiêu chí cuối khó chia thành nhãn cho từng
bước. Phần thưởng có thể là độ hoàn thành mục tiêu hoặc đánh giá chất
lượng. Tuy nhiên, một reward dễ tối ưu có thể khác mục tiêu con người
thực sự muốn.

Cần giữ bối cảnh thời gian: bài công bố năm 2020 và không phải survey về
toàn bộ các kỹ thuật căn chỉnh LLM xuất hiện sau đó. Báo cáo dùng nó cho
nền tảng RL và không gán các mô hình mới vào danh sách đóng góp nguyên
bản. So sánh với paper MAS năm 2025 cần nêu khoảng cách này.

\subsubsection{Autonomous IoT và computer
vision}\label{autonomous-iot-vuxe0-computer-vision}

Trong IoT, bài toán thường liên quan năng lượng, kết nối, thu thập dữ
liệu và phân bổ tài nguyên. Hành động có thể làm thay đổi thời lượng pin
hoặc mức thông tin nhận được. Với computer vision, RL có thể tham gia
lựa chọn vùng quan sát hoặc quyết định tuần tự, trong khi nhiều bài nhận
dạng thuần túy vẫn thích hợp hơn với supervised learning.

Bài học lựa chọn phương pháp: có ảnh hoặc có sensor chưa đủ để kết luận
cần RL. Nếu nhiệm vụ chỉ dự đoán nhãn từ một ảnh và hành động không ảnh
hưởng dữ liệu tương lai, classifier có thể phù hợp hơn. RL được biện
minh bởi cấu trúc quyết định và phản hồi, không phải loại dữ liệu đầu
vào.

\subsubsection{Games, self-organized systems và
scheduling}\label{games-self-organized-systems-vuxe0-scheduling}

Trò chơi cung cấp môi trường có luật rõ, vòng tương tác lặp và công cụ
đo reward, vì vậy thuận tiện cho nghiên cứu. Tuy nhiên, thành công trong
trò chơi có thông tin sạch không đồng nghĩa thành công trong hệ kinh
doanh có nhu cầu thay đổi và hạn chế can thiệp. Cần kiểm tra cơ chế
chuyển từ benchmark sang miền ứng dụng.

Trong scheduling, quyết định hiện tại có thể chiếm tài nguyên và ảnh
hưởng các công việc về sau. Ràng buộc cứng như deadline hoặc năng lực
máy cần được thể hiện trực tiếp. Nếu chỉ thêm một penalty nhỏ vào
reward, agent có thể chọn vi phạm để đổi lợi ích khác. Điều này tương tự
việc phải lọc sản phẩm vượt ngân sách trước khi xếp hạng trong tutorial.

\subsection{Các thách thức lớn và đọc phản
biện}\label{cuxe1c-thuxe1ch-thux1ee9c-lux1edbn-vuxe0-ux111ux1ecdc-phux1ea3n-biux1ec7n}

\subsubsection{Sample efficiency và độ bao
phủ}\label{sample-efficiency-vuxe0-ux111ux1ed9-bao-phux1ee7}

Paper nhấn mạnh chi phí dữ liệu và không gian trạng thái lớn. Đối với
lab, 5.000 episode tương đương 100.000 bước mỗi lượt training, nhưng
không có nghĩa mọi state--action được quan sát đủ cân bằng. Một trạng
thái chứa đối thủ giảm giá có thể hiếm khi policy tiến về không giảm
giá. Chỉ báo tổng số bước che khuất sự thiếu bao phủ cục bộ.

Đánh giá sample efficiency nên xem hiệu năng theo số tương tác và dùng
cùng ngân sách giữa thuật toán. So sánh một agent được train gấp mười
lần với baseline ít dữ liệu sẽ không tách được chất lượng phương pháp.
Một đề xuất mở rộng là lưu visit count cho từng state--action và phân
tích sai lệch Q ở các ô hiếm.

\subsubsection{Reward design và reward
hacking}\label{reward-design-vuxe0-reward-hacking}

Mục tiêu được số hóa có thể bỏ sót điều người dùng quan tâm. Nếu reward
là số đơn, agent dễ ưu tiên giảm giá mạnh dù lợi nhuận thấp. Nếu reward
là doanh thu, nó có thể bán hàng lỗ. Nếu reward chỉ là lợi nhuận ngắn
hạn, nó có thể không tính trải nghiệm khách hàng hoặc biến động giá.
Việc định nghĩa reward là một phần nội dung khoa học, không chỉ tham số
code.

Cải tiến phù hợp là tách objective và constraints: tối đa lợi nhuận
trong miền hành động tuân thủ các điều kiện về margin và ngân sách. Có
thể đưa một số ưu tiên mềm vào reward, nhưng phải công bố trọng số và
phân tích độ nhạy. Một con số tổng hợp cao không chứng minh mọi mặt đều
được cải thiện.

\subsubsection{Partial observability và
non-stationarity}\label{partial-observability-vuxe0-non-stationarity}

Khi chỉ quan sát một phần, agent có thể cần lịch sử hoặc belief state.
Trong lab, rule opponent dùng hành động trước của agent A, nhưng bảng Q
của A chỉ nhận nhu cầu và hành động trước của B. Hai lịch sử khác nhau
có cùng quan sát A nhưng dẫn đến hành động hiện tại khác nhau của B. Đây
là ví dụ cụ thể về state aliasing, ngay cả khi đối thủ không học.

Khi cả hai agent cùng học, chính sách đối thủ thay đổi qua episode nên
reward và chuyển tiếp hiệu dụng đối với một learner thay đổi. Frozen
evaluation sau training giúp đo chất lượng cặp policy tại một thời điểm,
nhưng không chứng minh quá trình học đã tiến đến cân bằng. Cần thêm
cross-play, nhiều seed và phân tích dao động hoặc best response.

\subsubsection{Bias--variance và những phát biểu cần điều
kiện}\label{biasvariance-vuxe0-nhux1eefng-phuxe1t-biux1ec3u-cux1ea7n-ux111iux1ec1u-kiux1ec7n}

Phần Challenges của paper so sánh MC, TD và biến động ước lượng. Báo cáo
giữ trực giác rằng target dài có thể nhiều nhiễu hơn, còn bootstrap có
thể mang sai lệch, nhưng không biến trực giác đó thành kết luận tuyệt
đối. Một estimator cụ thể có thể có hành vi khác tùy dữ liệu và cấu trúc
môi trường.

Paper cũng có những lời khẳng định rộng về dữ liệu tổng hợp và giảm
bias. Dữ liệu tổng hợp không bảo đảm loại thiên lệch; nó có thể tái tạo
hoặc làm nặng thiên lệch mô hình sinh. Trong tutorial, classifier đạt
100\% trên test tổng hợp nhưng vẫn nhận sai hai câu ngoài miền. Đây là
minh chứng thực tế vì sao cần tách chất lượng trên phân phối tạo dữ liệu
khỏi độ bền ngoài phân phối.

\subsection{Kiểm tra giả định trên chính môi trường
assignment}\label{kiux1ec3m-tra-giux1ea3-ux111ux1ecbnh-truxean-chuxednh-muxf4i-trux1b0ux1eddng-assignment}

Gọi trạng thái đầy đủ \(x_t=(d_t,u_t,v_t)\) với \(u_t=a^A_{t-1}\) và
\(v_t=a^B_{t-1}\). Với rule opponent, \(b_t=d_tu_t\). Quan sát tham
chiếu của A là \(o_t=(d_t,v_t)\), không có \(u_t\). Xét \(d_t=1,v_t=0\):
nếu \(u_t=0\) thì B chọn 0; nếu \(u_t=1\) thì B chọn 1. Cùng một quan
sát và cùng hành động A dẫn đến reward kỳ vọng khác nhau.

Ví dụ A không giảm: ở nhu cầu cao, B không giảm cho conversion 0.55 và
lợi nhuận kỳ vọng \(20\times0.55\times5=55\); B giảm cho conversion 0.45
và kỳ vọng 45. Chênh 10 không phải nhiễu ngẫu nhiên, mà do thông tin
thiếu. Vì A biết hành động mình vừa làm, việc bổ sung bit này là khả thi
và chi phí nhỏ.

Một yếu tố khác là thời điểm. Ở bước cuối, \(Q\) chỉ gồm reward ngay; ở
bước đầu còn nhiều reward tương lai. Nếu dùng một bảng cho cả hai mà
không có thời gian, cùng state--action nhận target theo những chân trời
khác nhau. Báo cáo giữ baseline tám giá trị để đúng tinh thần đề, rồi
dùng DP với đầy đủ trạng thái và thời gian để kiểm tra.

\subsection{Phân tích lợi nhuận kỳ vọng và chính sách
mốc}\label{phuxe2n-tuxedch-lux1ee3i-nhuux1eadn-kux1ef3-vux1ecdng-vuxe0-chuxednh-suxe1ch-mux1ed1c}

Với \(a,b\in\{0,1\}\), conversion ở các tham số mặc định luôn nằm giữa
0.20 và 0.70 nên clip không hoạt động. Reward kỳ vọng là:

\[\bar r(d,a,b)=(10+10d)(0.30+0.25d+0.15a-0.10b)(5-2a).\]

Nếu \(b=0\), nhu cầu thấp: không giảm thu 15, giảm thu 13.5; nhu cầu
cao: không giảm 55, giảm 42. Nếu \(b=1\), các giá trị lần lượt là 10 và
10.5 ở thấp, 45 và 36 ở cao. Vì rule opponent chỉ giảm khi nhu cầu cao
và A vừa giảm ở bước trước, trường hợp thấp mà B giảm không xuất hiện
trong baseline.

Không giảm là tốt hơn tức thời ở các trạng thái mà rule opponent tạo ra.
Nó còn tránh kích hoạt phản ứng giảm giá của B ở bước kế tiếp khi nhu
cầu cao. Quy hoạch động xác nhận hành động 0 ở mọi thời điểm và trạng
thái đầy đủ trong mô hình tham chiếu. Đây là best response với đối thủ
đã định nghĩa, không phải chứng minh luôn không giảm là tối ưu cho mọi
thị trường.

Khi cả hai khởi đầu với lịch sử 0 và dùng rule, không bên nào tạo ra tín
hiệu giảm giá nên cả hai luôn chọn 0. Nhu cầu thấp/cao có xác suất 1/2,
vì vậy lợi nhuận kỳ vọng mỗi bước là \((15+55)/2=35\). Kết quả đánh giá
thực nghiệm dao động quanh 35 do nhu cầu và sales; kỳ vọng 35 không được
chép thành mọi số đo của bảng kết quả.

\subsection{Cầu nối với MAS và phương pháp nghiên cứu
chung}\label{cux1ea7u-nux1ed1i-vux1edbi-mas-vuxe0-phux1b0ux1a1ng-phuxe1p-nghiuxean-cux1ee9u-chung}

Paper MAS nhấn mạnh kênh phối hợp, còn paper RL nhấn mạnh điều chỉnh
chính sách bằng feedback. Hai quan điểm bổ sung: một hệ có thể thiết kế
cơ chế giao tiếp rất rõ nhưng policy quyết định kém; một learner có
reward tốt trong mô phỏng vẫn có thể bị tích hợp sai vào một workflow.
Cần kiểm tra cả vòng học và hợp đồng vận hành.

Ví dụ nếu chuyển learner bán hàng thành một thành phần của hệ thương
mại, nó chỉ nên đề xuất mức giảm trong tập được phép. Agent chính sách
kiểm tra margin, inventory xác nhận khả năng giao hàng và hệ giao dịch
thực hiện với xác nhận tương ứng. Reward training không tự cấp quyền
thay đổi giá hoặc bỏ qua dữ liệu kho.

Một khuôn nghiên cứu chung là: mô hình hóa giả định; xây baseline đơn
giản; lưu dữ liệu thô; đánh giá đóng băng; phân tích theo loại trạng
thái; thử thay đổi môi trường; và nêu giới hạn suy rộng. Cách làm này có
giá trị hơn kết luận ``AI học được nên thông minh hơn luật'' từ một lượt
chạy.

\subsection{Kết luận riêng cho Paper
2}\label{kux1ebft-luux1eadn-riuxeang-cho-paper-2}

Paper cung cấp đường đi từ MDP đến các họ thuật toán và ứng dụng. Phần
có giá trị cho bài học là những câu hỏi nền tảng: trạng thái có đủ
không, reward đo gì, khám phá tốn gì, target được tính thế nào và kết
quả có được đánh giá độc lập không? Khi áp dụng vào assignment, các câu
hỏi này phát hiện những vấn đề cụ thể mà việc chỉ chạy chương trình tham
chiếu sẽ bỏ qua.

Báo cáo bổ sung ký hiệu kỳ vọng rõ, phân biệt Markov với tính dừng, hữu
hạn chân trời với continuing, hiệu năng training với evaluation và tối
ưu với hội tụ. Các thực nghiệm đi kèm không nhằm thay thế survey, mà
biến kiến thức tổng quan thành một trường hợp có thể tính tay, chạy lại
và phản biện bằng dữ liệu.
